
\subsubsection{6.1 Findings}

\textbf{Counterfactual density explains execution advantage.} The 84.7\% CF-score rate demonstrates that execution traces contain dense, extractable localization information. Type bugs (0.665) and logic bugs (0.583) provide the richest signals.

\textbf{The mechanism appears causal, not correlational.} The sequential experiments establish that localization enables targeting, which enables higher fix rates. Removing localization (binary feedback) breaks this chain despite providing ground truth about correctness.

\textbf{Complexity effect is inverted.} Contrary to the prediction that complex tasks would benefit more from precise localization, execution advantage is uniform or slightly higher on simpler tasks. This suggests a task difficulty ceiling: when problems require understanding multi-step dependencies, localization of individual errors provides diminishing returns.

\subsubsection{6.2 Limitations}

\textbf{L1: Partial Benchmark Coverage.} H-E1 processed 34/164 HumanEval problems due to GPU time constraints. H-M2 validation used only 5 problems. Effect size estimates may shift with full data.

\textbf{L2: Single Model Family.} All experiments used CodeLlama-7B-Instruct. Results may be model-specific.

\textbf{L3: Variable State Extraction Failed.} 0\% of traces had variable state extracted, despite parser implementation. CF-scores are conservative lower bounds.

\textbf{L4: AI Critic Model.} Using CodeLlama-7B as critic may underestimate stronger critics (GPT-4, Claude).

\subsubsection{6.3 Implications}

The findings suggest that practitioners designing code refinement pipelines should prioritize detailed error traces over simple test outcomes. The marginal cost of capturing line numbers and error types yields refinement improvement. A well-structured AI critic providing accurate localization could, in principle, approach execution feedback performance---the key is information structure, not information source.

For complex, distributed bugs, complementary strategies (e.g., divide-and-conquer, hierarchical debugging) may be needed alongside localization.
